\providecommand{\ourmethod}{\textsc{SkillOpt}}
\providecommand{\na}{--}
\providecommand{\harnessrow}[1]{\rowcolor{black!7}\multicolumn{8}{l}{\textbf{#1}}}
% Blue tint applied per cell. The \dvalp / \dvaln macros used in the data
% cells embed an invisible left phantom of the same width as the right
% subscript so the black number stays centered in the column AND the
% colored subscript lives entirely inside this cell (no clipping by the
% next cell's background).
\newcommand{\methodcell}[1]{\cellcolor{blue!8}#1}
\newcommand{\ablationbenchmarkheader}{\textbf{Setting} & \textbf{SearchQA} & \textbf{Spreadsheet} & \textbf{LiveMath} \\}

\section{Method}
\label{sec:method}

\begin{figure}[t]
    \centering
    \includegraphics[width=\textwidth]{sections/figures/pipeline.pdf}
    \caption{Pipeline of \ourmethod{}. A frozen target model executes a rollout batch with the current skill; an optimizer model performs minibatch reflection over successes and failures, proposes bounded add/delete/replace edits, merges and ranks them under a scheduled edit budget, and accepts the candidate skill only through a held-out validation gate. Across epochs, the slow/meta update retains longer-horizon lessons without changing the target model.}
    \label{fig:skillopt_pipeline}
\end{figure}

% \ourmethod{} optimizes an agent skill as the external trainable state for adapting a frozen target model, building on language-level optimization and reusable-skill views of agent adaptation \citep{zhou2023large,yang2024large,khattab2024dspy,li2026skillsbench,jiang2026sok}. The central advantage is that domain adaptation happens in a compact text artifact rather than in model weights: the target model, backend, and execution harness remain fixed, while the skill document learns reusable procedures for tool use, evidence gathering, verifier-facing formatting, and domain-specific failure avoidance. Figure~\ref{fig:skillopt_pipeline} shows the full loop. It is deliberately shaped like a training algorithm: a forward pass collects rollout evidence, a language-level backward pass converts that evidence into edit directions, and optimizer controls decide how far the skill is allowed to move.

\subsection{Problem Setup}
\label{sec:method_skill}

A skill $s$ is a natural-language policy inserted into the agent context before execution, consistent with recent work treating skills as reusable procedural knowledge for agents \citep{li2026skillsbench,jiang2026sok}. In direct-chat benchmarks, it is prepended to the system or developer instruction; in tool-use harnesses, it becomes persistent procedural memory. We use $M$ to denote the frozen target model whose behavior is being adapted through skill optimization. For a harness $h$, task $x$, and skill $s$, execution produces a trajectory $\tau$ and a scalar score $r$:
\begin{equation}
    (\tau(s), r(s)) = h(M, x, s), \qquad r(s) \in [0,1].
\end{equation}

Given train, selection, and test splits $D_{\mathrm{tr}},D_{\mathrm{sel}},D_{\mathrm{test}}$, \ourmethod{} uses $D_{\mathrm{tr}}$ to generate a set of candidate skills $\mathcal{C}(D_{\mathrm{tr}})$, selects the best skill on $D_{\mathrm{sel}}$, and reports the final performance on $D_{\mathrm{test}}$:
\begin{equation}
    s^\star_{\mathrm{sel}}
    =
    \arg\max_{s \in \mathcal{C}(D_{\mathrm{tr}})}
    \frac{1}{|D_{\mathrm{sel}}|}
    \sum_{x \in D_{\mathrm{sel}}}
    r(s),
\end{equation}
\begin{equation}
    \mathrm{Test}(s^\star_{\mathrm{sel}})
    =
    \frac{1}{|D_{\mathrm{test}}|}
    \sum_{x \in D_{\mathrm{test}}}
    r(s^\star_{\mathrm{sel}}).
\end{equation}
The training split supplies experience, the selection split gates updates, and the test split is used only for final reporting. The optimizer state contains the current skill, the best validation-gated skill, cached skill hashes, an epoch-local rejected-step buffer, and optional slow/meta-update state. Only the best accepted skill is exported as \texttt{best\_skill.md}.

\subsection{Forward Pass: Rollout Evidence}
\label{sec:method_forward}

At each optimization step, the target model runs a rollout batch from $D_{\mathrm{tr}}$ with the current skill. The harness records task metadata, messages, tool calls, observations, command outputs, final answers, verifier feedback, and benchmark-specific context such as spreadsheet previews, document references, or compact execution traces. This batch is the evidence unit: small batches update quickly but noisily, while larger batches expose more recurring patterns before the skill changes. The implementation also supports accumulation, where several rollout batches are reflected on separately and merged into one update, decoupling execution throughput from update frequency.

\subsection{Backward Pass: Minibatch Reflection}
\label{sec:method_backward}

The optimizer model turns trajectories into skill edits, following the broader line of trajectory-driven reflection and prompt evolution \citep{shinn2023reflexion,madaan2023selfrefine,agrawal2025gepa}. It first separates failures from successes and partitions each group into reflection minibatches. This matters because single trajectories often produce anecdotal fixes, while minibatches expose reusable procedural errors: the agent consistently searches the wrong source, writes an answer in the wrong format, or fails to verify a tool result. Failure minibatches propose missing or corrective rules; success minibatches preserve behaviors that already work. Each reflection returns structured add/delete/replace edits, or in rewrite mode a small set of rewrite suggestions.

Local proposals are merged hierarchically by first consolidating failure- and success-driven edits separately, then combining them with priority on failure corrections. This step filters duplicate, contradictory, and example-specific suggestions before the optimizer selects the final bounded update.

\subsection{Bounded Text Updates}
\label{sec:method_controls}

The learning-rate analogue in \ourmethod{} is the edit budget $L_t$: the maximum number of skill edits applied at step $t$. After aggregation, the optimizer model ranks the merged edit pool by expected utility and clips it to the top $L_t$ edits. This is the key difference from ad hoc prompt rewriting. Unbounded rewrites can erase useful rules, introduce incompatible instructions, or overfit to a local failure; bounded updates preserve continuity while still allowing the skill to acquire new procedures. \ourmethod{} supports constant, linear, cosine, and autonomous schedules. The default cosine schedule starts with larger edits and decays toward smaller consolidation steps.

The selected edits produce a candidate skill. In patch mode, edits are localized operations such as append, insert, replace, and delete; in rewrite mode, selected suggestions condition a full skill rewrite. Step-level edits cannot overwrite the protected slow-update field, so fast local changes and slower epoch-wise consolidation remain separated.


\subsection{Validation Gate and Rejected-Edit Buffer}
\label{sec:method_gate}

Every candidate skill is evaluated on $D_{\mathrm{sel}}$ with the same frozen target model and harness. If it improves over the current selection score, it becomes the new current skill; if it also exceeds the best score so far, it becomes \texttt{best\_skill.md}. Otherwise it is rejected. This gate turns reflection into propose-and-test optimization rather than unconditional self-editing, which is crucial because plausible textual diagnoses can still hurt the actual target model.

Rejected updates are still useful. The optimizer records an epoch-local buffer containing observed failure patterns and, for rejected steps, the edits that were tried and the score drop they caused. Later reflection calls in the same epoch receive this buffer, so the optimizer model can avoid repeating failed edits and focus on unresolved failures. This gives the loop negative feedback during training without adding inference-time cost.
\begin{table*}[!htbp]
\centering
\scriptsize
\setlength{\tabcolsep}{1.5pt}
\renewcommand{\arraystretch}{1.0}
\begin{tabular}{l l c c c c c c}
\toprule
\textbf{Model} & \textbf{Skill source} & \textbf{SearchQA} & \textbf{Spreadsheet} & \textbf{OfficeQA} & \textbf{DocVQA} & \textbf{LiveMath} & \textbf{ALFWorld} \\
\midrule
\harnessrow{No harness / direct chat} \\
\multirow{7}{*}{GPT--5.5} & No skill & 77.7 & 41.8 & 33.1 & 78.8 & 37.6 & 83.6 \\
 & Human skill & \dvalp{81.8}{4.1} & \dvalp{72.9}{31.1} & \dvalp{\underline{66.9}}{33.8} & \dvalp{90.1}{11.3} & \dvalp{38.4}{0.8} & \dvalp{91.8}{8.2} \\
 & LLM skill & \dvalp{80.9}{3.2} & \dvalp{43.2}{1.4} & \dvalp{51.7}{18.6} & \dvalp{89.6}{10.8} & \dvalp{40.0}{2.4} & \dvalp{\underline{93.3}}{9.7} \\
 & Trace2Skill & \dvalp{82.4}{4.7} & \dvalp{49.6}{7.8} & \dvalp{65.7}{32.6} & \dvalp{\underline{90.6}}{11.8} & \dvalp{\underline{52.0}}{14.4} & \dvalp{87.3}{3.7} \\
 & TextGrad & \dvalp{81.4}{3.7} & \dvaln{41.1}{0.7} & \dvalp{42.0}{8.9} & \dvalp{87.2}{8.4} & \dvalp{49.2}{11.6} & \dvaln{82.8}{0.8} \\
 & GEPA & \dvalp{\underline{84.8}}{7.1} & \dvalp{\underline{73.6}}{31.8} & \dvalp{63.9}{30.8} & \dvalp{89.1}{10.3} & \dvalp{43.2}{5.6} & \dvalp{85.8}{2.2} \\
 & \methodcell{\textbf{\ourmethod{}}} & \methodcell{\dvalp{\textbf{87.3}}{9.6}} & \methodcell{\dvalp{\textbf{80.7}}{38.9}} & \methodcell{\dvalp{\textbf{72.1}}{39.0}} & \methodcell{\dvalp{\textbf{91.2}}{12.4}} & \methodcell{\dvalp{\textbf{66.9}}{29.3}} & \methodcell{\dvalp{\textbf{95.5}}{11.9}} \\
\cmidrule(lr){1-8}
\multirow{7}{*}{GPT--5.4} & No skill & 76.9 & 41.4 & 50.0 & 77.6 & 36.8 & 75.4 \\
 & Human skill & \dvalp{80.1}{3.2} & \dvalp{59.3}{17.9} & \dvalp{59.3}{9.3} & \dvalp{88.5}{10.9} & \dvalp{\underline{41.6}}{4.8} & \dvaln{74.6}{0.8} \\
 & LLM skill & \dvalp{79.4}{2.5} & \dvalp{46.1}{4.7} & \dvaln{20.4}{29.6} & \dvalp{\underline{90.4}}{12.8} & \dvalp{36.8}{0.0} & \dvalp{\underline{83.6}}{8.2} \\
 & Trace2Skill & \dvalp{81.2}{4.3} & \dvalp{53.2}{11.8} & \dvalp{51.2}{1.2} & \dvalp{89.3}{11.7} & \dvalp{40.0}{3.2} & \dvaln{73.1}{2.3} \\
 & TextGrad & \dvalp{\underline{82.5}}{5.6} & \dvaln{38.6}{2.8} & \dvalp{54.7}{4.7} & \dvalp{85.3}{7.7} & \dvaln{36.6}{0.2} & \dvalp{78.4}{3.0} \\
 & GEPA & \dvalp{82.4}{5.5} & \dvalp{\underline{61.1}}{19.7} & \dvalp{\underline{60.4}}{10.4} & \dvalp{88.3}{10.7} & \dvalp{\underline{41.6}}{4.8} & \dvaln{74.6}{0.8} \\
 & \methodcell{\textbf{\ourmethod{}}} & \methodcell{\dvalp{\textbf{83.1}}{6.2}} & \methodcell{\dvalp{\textbf{62.5}}{21.1}} & \methodcell{\dvalp{\textbf{62.8}}{12.8}} & \methodcell{\dvalp{\textbf{91.2}}{13.6}} & \methodcell{\dvalp{\textbf{44.0}}{7.2}} & \methodcell{\dvalp{\textbf{91.0}}{15.6}} \\
\cmidrule(lr){1-8}
\multirow{7}{*}{GPT--5.4-mini} & No skill & 75.9 & 36.1 & 22.1 & 71.4 & 14.7 & 73.1 \\
 & Human skill & \dvalp{77.2}{1.3} & \dvalp{\underline{42.9}}{6.8} & \dvalp{\underline{45.9}}{23.8} & \dvalp{85.0}{13.6} & \dvalp{\underline{28.8}}{14.1} & \dvaln{56.7}{16.4} \\
 & LLM skill & \dvalp{76.1}{0.2} & \dvalp{36.8}{0.7} & \dvalp{36.6}{14.5} & \dvalp{86.4}{15.0} & \dvalp{28.0}{13.3} & \dvaln{65.7}{7.4} \\
 & Trace2Skill & \dvalp{78.6}{2.7} & \dvalp{40.7}{4.6} & \dvaln{20.9}{1.2} & \dvalp{\underline{88.5}}{17.1} & \dvalp{\textbf{32.8}}{18.1} & \dvalp{\underline{82.8}}{9.7} \\
 & TextGrad & \dvalp{77.5}{1.6} & \dvalp{38.2}{2.1} & \dvalp{30.0}{7.9} & \dvalp{84.0}{12.6} & \dvalp{27.2}{12.5} & \dvaln{70.9}{2.2} \\
 & GEPA & \dvalp{\underline{79.4}}{3.5} & \dvalp{42.5}{6.4} & \dvalp{45.3}{23.2} & \dvalp{83.7}{12.3} & \dvalp{27.2}{12.5} & \dvalp{81.3}{8.2} \\
 & \methodcell{\textbf{\ourmethod{}}} & \methodcell{\dvalp{\textbf{80.2}}{4.3}} & \methodcell{\dvalp{\textbf{47.5}}{11.4}} & \methodcell{\dvalp{\textbf{48.8}}{26.7}} & \methodcell{\dvalp{\textbf{90.9}}{19.5}} & \methodcell{\dvalp{\textbf{32.8}}{18.1}} & \methodcell{\dvalp{\textbf{85.8}}{12.7}} \\
\cmidrule(lr){1-8}
\multirow{7}{*}{GPT--5.4-nano} & No skill & 55.8 & 23.5 & 16.3 & 30.8 & 23.2 & 34.3 \\
 & Human skill & \dvalp{66.1}{10.3} & \dvalp{\underline{41.8}}{18.3} & \dvalp{\underline{46.5}}{30.2} & \dvalp{73.5}{42.7} & \dvaln{16.8}{6.4} & \dvaln{29.9}{4.4} \\
 & LLM skill & \dvalp{62.4}{6.6} & \dvalp{33.6}{10.1} & \dvaln{12.8}{3.5} & \dvalp{71.7}{40.9} & \dvaln{20.0}{3.2} & \dvalp{65.7}{31.4} \\
 & Trace2Skill & \dvalp{69.9}{14.1} & \dvalp{35.4}{11.9} & \dvalp{16.3}{0.0} & \dvalp{\underline{77.3}}{46.5} & \dvalp{\underline{25.6}}{2.4} & \dvalp{58.2}{23.9} \\
 & TextGrad & \dvalp{\underline{73.4}}{17.6} & \dvalp{32.5}{9.0} & \dvalp{38.7}{22.4} & \dvalp{68.3}{37.5} & \dvaln{20.8}{2.4} & \dvalp{42.5}{8.2} \\
 & GEPA & \dvalp{73.2}{17.4} & \dvalp{37.2}{13.7} & \dvalp{45.0}{28.7} & \dvalp{71.2}{40.4} & \dvalp{24.8}{1.6} & \dvalp{\underline{68.7}}{34.4} \\
 & \methodcell{\textbf{\ourmethod{}}} & \methodcell{\dvalp{\textbf{74.8}}{19.0}} & \methodcell{\dvalp{\textbf{42.5}}{19.0}} & \methodcell{\dvalp{\textbf{50.0}}{33.7}} & \methodcell{\dvalp{\textbf{80.2}}{49.4}} & \methodcell{\dvalp{\textbf{27.2}}{4.0}} & \methodcell{\dvalp{\textbf{69.4}}{35.1}} \\
\cmidrule(lr){1-8}
\multirow{7}{*}{GPT--5.2} & No skill & 71.9 & 38.2 & 34.9 & 73.1 & 20.8 & 68.7 \\
 & Human skill & \dvalp{75.8}{3.9} & \dvalp{48.2}{10.0} & \dvalp{\underline{55.2}}{20.3} & \dvalp{89.0}{15.9} & \dvalp{23.2}{2.4} & \dvaln{56.7}{12.0} \\
 & LLM skill & \dvalp{73.7}{1.8} & \dvalp{38.6}{0.4} & \dvaln{14.0}{20.9} & \dvalp{87.7}{14.6} & \dvalp{25.6}{4.8} & \dvalp{68.7}{0.0} \\
 & Trace2Skill & \dvalp{79.2}{7.3} & \dvalp{43.9}{5.7} & \dvalp{43.0}{8.1} & \dvalp{87.7}{14.6} & \dvalp{\underline{28.8}}{8.0} & \dvalp{76.9}{8.2} \\
 & TextGrad & \dvalp{82.6}{10.7} & \dvaln{32.1}{6.1} & \dvalp{51.1}{16.2} & \dvalp{87.2}{14.1} & \dvalp{22.4}{1.6} & \dvalp{76.1}{7.4} \\
 & GEPA & \dvalp{\underline{82.7}}{10.8} & \dvalp{\underline{54.3}}{16.1} & \dvalp{53.4}{18.5} & \dvalp{\underline{89.5}}{16.4} & \dvalp{\underline{28.8}}{8.0} & \dvalp{\underline{82.8}}{14.1} \\
 & \methodcell{\textbf{\ourmethod{}}} & \methodcell{\dvalp{\textbf{83.1}}{11.2}} & \methodcell{\dvalp{\textbf{57.1}}{18.9}} & \methodcell{\dvalp{\textbf{56.4}}{21.5}} & \methodcell{\dvalp{\textbf{89.6}}{16.5}} & \methodcell{\dvalp{\textbf{36.0}}{15.2}} & \methodcell{\dvalp{\textbf{85.1}}{16.4}} \\
\cmidrule(lr){1-8}
\multirow{7}{*}{Qwen3.5--4B} & No skill & 68.1 & 9.3 & 14.5 & 86.9 & 22.4 & 30.6 \\
 & Human skill & \dvaln{66.3}{1.8} & \dvalp{16.4}{7.1} & \dvalp{\underline{22.7}}{8.2} & \dvalp{87.8}{0.9} & \dvaln{18.4}{4.0} & \dvaln{28.4}{2.2} \\
 & LLM skill & \dvaln{65.0}{3.1} & \dvalp{10.4}{1.1} & \dvalp{20.9}{6.4} & \dvalp{\underline{88.0}}{1.1} & \dvalp{\underline{28.8}}{6.4} & \dvalp{55.2}{24.6} \\
 & Trace2Skill & \dvalp{68.5}{0.4} & \dvalp{\underline{19.3}}{10.0} & \dvalp{16.3}{1.8} & \dvalp{\underline{88.0}}{1.1} & \dvalp{27.2}{4.8} & \dvalp{\underline{64.9}}{34.3} \\
 & TextGrad & \dvaln{60.7}{7.4} & \dvalp{13.9}{4.6} & \dvalp{20.9}{6.4} & \dvaln{85.6}{1.3} & \dvaln{10.6}{11.8} & \dvalp{53.7}{23.1} \\
 & GEPA & \dvalp{\underline{68.6}}{0.5} & \dvalp{16.9}{7.6} & \dvalp{\underline{22.7}}{8.2} & \dvaln{85.1}{1.8} & \dvalp{\underline{28.8}}{6.4} & \dvalp{60.4}{29.8} \\
 & \methodcell{\textbf{\ourmethod{}}} & \methodcell{\dvalp{\textbf{71.2}}{3.1}} & \methodcell{\dvalp{\textbf{23.9}}{14.6}} & \methodcell{\dvalp{\textbf{29.7}}{15.2}} & \methodcell{\dvalp{\textbf{89.0}}{2.1}} & \methodcell{\dvalp{\textbf{52.0}}{29.6}} & \methodcell{\dvalp{\textbf{81.3}}{50.7}} \\
\cmidrule(lr){1-8}
\multirow{7}{*}{Qwen3.6--35B-A3B} & No skill & 72.7 & 38.2 & 45.9 & 87.6 & \underline{31.2} & 59.7 \\
 & Human skill & \dvalp{74.1}{1.4} & \dvalp{44.3}{6.1} & \dvaln{41.9}{4.0} & \dvalp{88.4}{0.8} & \dvaln{29.6}{1.6} & \dvaln{44.8}{14.9} \\
 & LLM skill & \dvaln{72.6}{0.1} & \dvalp{42.9}{4.7} & \dvalp{\underline{46.5}}{0.6} & \dvalp{88.4}{0.8} & \dvaln{24.8}{6.4} & \dvalp{60.4}{0.7} \\
 & Trace2Skill & \dvalp{75.4}{2.7} & \dvaln{33.2}{5.0} & \dvaln{32.0}{13.9} & \dvalp{\underline{90.4}}{2.8} & \dvaln{29.6}{1.6} & \dvalp{70.9}{11.2} \\
 & TextGrad & \dvalp{\underline{76.4}}{3.7} & \dvaln{22.9}{15.3} & \dvaln{33.7}{12.2} & \dvaln{84.5}{3.1} & \dvaln{7.2}{24.0} & \dvalp{67.9}{8.2} \\
 & GEPA & \dvalp{75.8}{3.1} & \dvalp{\underline{45.4}}{7.2} & \dvaln{43.6}{2.3} & \dvalp{88.0}{0.4} & \dvalp{\underline{31.2}}{0.0} & \dvalp{\underline{73.9}}{14.2} \\
 & \methodcell{\textbf{\ourmethod{}}} & \methodcell{\dvalp{\textbf{80.3}}{7.6}} & \methodcell{\dvalp{\textbf{47.5}}{9.3}} & \methodcell{\dvalp{\textbf{47.1}}{1.2}} & \methodcell{\dvalp{\textbf{91.4}}{3.8}} & \methodcell{\dvalp{\textbf{41.6}}{10.4}} & \methodcell{\dvalp{\textbf{82.1}}{22.4}} \\
\midrule
\harnessrow{Codex harness} \\
\multirow{5}{*}{GPT--5.5} & No skill & 81.8 & 27.5 & 38.3 & 87.2 & 35.2 & \na \\
 & Human skill & \dvalp{\underline{84.1}}{2.3} & \dvalp{50.7}{23.2} & \dvalp{40.0}{1.7} & \dvalp{88.8}{1.6} & \dvalp{48.8}{13.6} & \na \\
 & LLM skill & \dvalp{83.4}{1.6} & \dvaln{25.0}{2.5} & \dvaln{34.3}{4.0} & \dvalp{\underline{89.8}}{2.6} & \dvalp{45.6}{10.4} & \na \\
 & EvoSkill & \dvaln{61.4}{20.4} & \dvalp{\underline{67.5}}{40.0} & \dvalp{\underline{42.4}}{4.1} & \dvalp{89.3}{2.1} & \dvalp{\underline{63.2}}{28.0} & \na \\
 & \methodcell{\textbf{\ourmethod{}}} & \methodcell{\dvalp{\textbf{87.3}}{5.5}} & \methodcell{\dvalp{\textbf{85.0}}{57.5}} & \methodcell{\dvalp{\textbf{51.1}}{12.8}} & \methodcell{\dvalp{\textbf{92.2}}{5.0}} & \methodcell{\dvalp{\textbf{78.4}}{43.2}} & \methodcell{\na} \\
\midrule
\harnessrow{Claude Code harness} \\
\multirow{5}{*}{GPT--5.5} & No skill & 81.9 & 22.1 & 57.6 & 86.6 & 40.8 & \na \\
 & Human skill & \dvalp{83.7}{1.8} & \dvalp{37.1}{15.0} & \dvalp{66.3}{8.7} & \dvalp{88.0}{1.4} & \dvalp{44.0}{3.2} & \na \\
 & LLM skill & \dvalp{82.4}{0.5} & \dvalp{37.1}{15.0} & \dvaln{56.4}{1.2} & \dvalp{\underline{89.6}}{3.0} & \dvalp{44.8}{4.0} & \na \\
 & EvoSkill & \dvalp{\underline{84.0}}{2.1} & \dvalp{\underline{75.0}}{52.9} & \dvalp{\underline{70.3}}{12.7} & \dvalp{87.2}{0.6} & \dvalp{\underline{52.0}}{11.2} & \na \\
 & \methodcell{\textbf{\ourmethod{}}} & \methodcell{\dvalp{\textbf{85.9}}{4.0}} & \methodcell{\dvalp{\textbf{80.4}}{58.3}} & \methodcell{\dvalp{\textbf{71.5}}{13.9}} & \methodcell{\dvalp{\textbf{90.1}}{3.5}} & \methodcell{\dvalp{\textbf{56.5}}{15.7}} & \methodcell{\na} \\
\bottomrule
\end{tabular}
\caption{Main results on held-out test splits. Scores are percentages; within each model--harness block, bold marks the best measured entry and underlining marks the second-best entry for each benchmark. Blue cells denote \ourmethod{}, and small green/red subscripts show the absolute change relative to the \emph{No skill} row of the same model in the same harness. We omit ALFWorld under Codex and Claude Code harnesses because ALFWorld requires persistent embodied-environment interaction. \ourmethod{} is the best-or-tied entry on every measured cell of the table, with positive gains over the no-skill baseline throughout.}
\label{tab:main_results_by_harness}
\end{table*}

\subsection{Epoch-Wise Slow/Meta Update}
\label{sec:method_slow_meta}

Fast updates learn from the current batch; the epoch-wise slow/meta update learns from adjacent epochs. At the end of an epoch, \ourmethod{} samples the same training items under the previous epoch's skill and the current skill, then groups them into improvements, regressions, persistent failures, and stable successes. The optimizer model writes a concise longitudinal guidance block into a protected slow-update field, and this candidate is still passed through the validation gate. Thus slow update captures durable domain lessons while preserving the same safety check as step-level edits.

The meta skill is optimizer-side only. It summarizes which edit patterns helped, which were rejected, and which failures persisted across epochs. This meta guidance is prepended to future optimizer prompts for reflection, merging, and ranking, but it is not shipped with the target model. The advantage is separation of concerns: the deployed skill remains compact and portable, while training benefits from a richer record of the editing process.

\subsection{Harness-Agnostic Deployment}
\label{sec:method_harness}

\ourmethod{} is harness-agnostic through a lightweight adapter interface, matching the broader trend toward agents embedded in tool-use and software-execution environments \citep{yao2022react,schick2023toolformer,yang2024sweagent}. An adapter constructs train/evaluation batches, injects the current skill into the agent context, runs the native harness, and returns scored trajectories. The same optimizer therefore works for direct QA, spreadsheet execution, document reasoning, multimodal QA, embodied environments, and Codex-style or Claude Code-style execution loops. This is the main practical advantage of treating skills as the adaptation layer: a stronger optimizer model can train a reusable skill artifact offline, and the resulting \texttt{best\_skill.md} can then be deployed or tested across target models, harnesses, and nearby benchmarks without changing model weights.

\begin{table*}[t]
\centering
\scriptsize
\setlength{\tabcolsep}{1.2pt}
\renewcommand{\arraystretch}{0.74}
\resizebox{\textwidth}{!}{%
\begin{minipage}{1.12\textwidth}
\begin{minipage}[c]{0.33\textwidth}
\centering
\textbf{(a) Training set size}\\[-1pt]
\begin{tabular}{@{}lccc@{}}
\toprule
\ablationbenchmarkheader
\midrule
1 example & 81.0 & 47.5 & 59.1 \\
20\% train & 84.1 & 69.0 & 65.9 \\
40\% train & \textbf{86.1} & 73.5 & 64.8 \\
80\% train & \textbf{86.1} & 77.6 & 67.0 \\
100\% train & 84.1 & \textbf{78.0} & \textbf{70.5} \\
\bottomrule
\end{tabular}
\end{minipage}
\hfill
\begin{minipage}[c]{0.33\textwidth}
\centering
\textbf{(b) Mini-batchsize}\\[-1pt]
\begin{tabular}{@{}lccc@{}}
\toprule
\ablationbenchmarkheader
\midrule
1 & 85.9 & 75.4 & 60.5 \\
2 & 86.3 & 77.1 & 54.8 \\
4 & 86.9 & 75.4 & \textbf{64.5} \\
8 & \textbf{87.1} & 77.5 & 61.3 \\
16 & 87.0 & \textbf{77.9} & 61.3 \\
32 & 86.9 & 77.5 & 58.9 \\
\bottomrule
\end{tabular}
\end{minipage}
\hfill
\begin{minipage}[c]{0.33\textwidth}
\centering
\textbf{(c) Batchsize}\\[-1pt]
\begin{tabular}{@{}lccc@{}}
\toprule
\ablationbenchmarkheader
\midrule
8 & 85.1 & 76.8 & 58.1 \\
24 & 86.4 & 77.1 & \textbf{62.9} \\
40 & 87.1 & \textbf{77.5} & 61.3 \\
56 & 86.5 & 76.8 & 56.5 \\
full epoch & \textbf{87.2} & 75.0 & 53.2 \\
\bottomrule
\end{tabular}
\end{minipage}

\vspace{2pt}

\begin{minipage}[c]{0.33\textwidth}
\centering
\textbf{(d) Learning rate}\\[-1pt]
\begin{tabular}{@{}lccc@{}}
\toprule
\ablationbenchmarkheader
\midrule
lr=1 & 85.5 & 77.5 & 62.1 \\
lr=2 & 86.7 & 77.5 & 60.5 \\
lr=4 & 86.5 & \textbf{78.2} & 56.5 \\
lr=8 & \textbf{87.0} & 73.6 & \textbf{66.9} \\
lr=16 & 86.8 & \textbf{78.2} & 65.3 \\
\bottomrule
\end{tabular}
\end{minipage}
\hfill
\begin{minipage}[c]{0.33\textwidth}
\centering
\textbf{(e) Learning-rate scheduler}\\[-1pt]
\begin{tabular}{@{}lccc@{}}
\toprule
\ablationbenchmarkheader
\midrule
constant & \textbf{87.3} & \textbf{80.7} & 62.1 \\
cosine & 87.1 & 77.5 & 61.3 \\
linear & 87.2 & 72.9 & \textbf{62.9} \\
\bottomrule
\end{tabular}
\end{minipage}
\hfill
\begin{minipage}[c]{0.33\textwidth}
\centering
\textbf{(f) Slow-update samples}\\[-1pt]
\begin{tabular}{@{}lccc@{}}
\toprule
\ablationbenchmarkheader
\midrule
5 & 86.8 & 76.4 & 64.5 \\
10 & 86.4 & 74.3 & \textbf{65.3} \\
20 & \textbf{87.1} & \textbf{77.5} & 61.3 \\
40 & 86.9 & 75.4 & 54.8 \\
\bottomrule
\end{tabular}
\end{minipage}
\end{minipage}%
}
\caption{Hyperparameter analysis for the text optimizer. Each panel changes one scalar or scheduling factor from the default setting unless noted. Panel (a) fixes the split to $4{:}1{:}5$ train/selection/test; the 1-example, 20\%, 40\%, and 80\% rows use subsets of the training partition, and the 100\% row reuses the completed $4{:}1{:}5$ split-ratio run. Panel (b) sweeps the reflection mini-batchsize $B_m$; panel (c) sweeps the rollout batchsize $B$.}
\label{tab:ablation_sweeps}
\end{table*}

\begin{table*}[t]
\centering
\scriptsize
\setlength{\tabcolsep}{4pt}
\renewcommand{\arraystretch}{0.82}
\resizebox{0.92\textwidth}{!}{%
\begin{tabular}{l l c c c}
\toprule
\textbf{Component} & \textbf{Setting} & \textbf{SearchQA} & \textbf{SpreadsheetBench} & \textbf{LiveMath} \\
\midrule
\rowcolor{blue!8}Learning-rate form & lr=4 (default) & \textbf{87.1} & \textbf{77.5} & \textbf{61.3} \\
 & dynamic lr & 85.8 & 71.8 & 54.0 \\
 & without lr & 84.6 & 75.7 & 57.3\\
\cmidrule(lr){1-5}
\rowcolor{blue!8}Rejected buffer & with rejected buffer & \textbf{87.1} & \textbf{77.5} & \textbf{61.3} \\
 & without rejected buffer & 85.5 & 72.9 & 58.9 \\
\cmidrule(lr){1-5}
\rowcolor{blue!8}Slow/meta update & meta skill and slow update & \textbf{87.1} & \textbf{77.5} & \textbf{61.3} \\
 & without meta skill & 85.1 & 75.7 & 58.1 \\
 & without meta skill and slow update & 86.3 & 55.0 & 59.7 \\
\bottomrule
\end{tabular}%
}
\caption{Component ablations for learning-rate form, rejected buffer, and epoch-wise slow/meta update. Light-blue rows mark the default setting within each component group; the learning-rate group uses the default lr=4 setting. Bold values mark the best measured result within that group and benchmark. The without-rejected-buffer row uses the matched no-buffer ablation setting.}
\label{tab:component_ablation}
\end{table*}

% =============================================================
% Merged transfer table (cross-model / cross-harness / cross-benchmark).
% The three legacy tables below this block are commented out for archival;
% all three legacy labels (tab:transfer_cross_model, tab:transfer_cross_harness,
% tab:transfer_cross_benchmark) are re-issued inside the merged table so that
% existing \ref{...} calls keep resolving to the right section.
% =============================================================
\begin{table*}[t]
\centering
\scriptsize
\setlength{\tabcolsep}{6pt}
\renewcommand{\arraystretch}{0.9}
\resizebox{0.95\textwidth}{!}{%
\begin{tabular}{l l l c c c}
\toprule
% --- (a) Cross-model transfer ---
\rowcolor{black!7}\multicolumn{6}{l}{\textbf{(a) Cross-model transfer}} \\
\midrule
\textbf{Source model} & \textbf{Target model} & \textbf{Benchmark} & \textbf{Baseline} & \textbf{Direct} & \textbf{Transferred} \\
\midrule
\multirow{3}{*}{GPT--5.4} & GPT--5.4      & \multirow{3}{*}{SpreadsheetBench} & 41.4 & 62.5 & \na \\
                          & GPT--5.4-mini &                                    & 36.1 & 47.5 & \dvalp{45.5}{9.4}  \\
                          & GPT--5.4-nano &                                    & 23.5 & 42.5 & \dvalp{26.5}{3.0}  \\
\cmidrule(lr){1-6}
\multirow{3}{*}{GPT--5.4} & GPT--5.4      & \multirow{3}{*}{LiveMath}         & 36.8 & 44.0 & \na \\
                          & GPT--5.4-mini &                                    & 14.7 & 32.8 & \dvalp{19.2}{4.5}  \\
                          & GPT--5.4-nano &                                    & 23.2 & 27.2 & \dvalp{28.8}{5.6}  \\
% --- (b) Cross-harness transfer ---
\midrule
\rowcolor{black!7}\multicolumn{6}{l}{\textbf{(b) Cross-harness transfer}} \\
\midrule
\textbf{Source harness} & \textbf{Target harness} & \textbf{Benchmark} & \textbf{Baseline} & \textbf{Direct} & \textbf{Transferred} \\
\midrule
Codex                       & Claude Code & \multirow{2}{*}{LiveMath}         & 40.8 & 56.5 & \dvalp{42.4}{1.6}  \\
Claude Code                 & Codex       &                                    & 35.2 & 78.4 & \dvalp{48.0}{12.8} \\
Codex                       & Claude Code & \multirow{2}{*}{SpreadsheetBench} & 22.1 & 80.4 & \dvalp{81.8}{59.7} \\
Claude Code                 & Codex       &                                    & 27.5 & 85.0 & \dvalp{71.1}{43.6} \\
% --- (c) Cross-benchmark transfer (math); LiveMath-as-source removed. ---
\midrule
\rowcolor{black!7}\multicolumn{6}{l}{\textbf{(c) Cross-benchmark transfer}} \\
\midrule
\textbf{Source benchmark} & \textbf{Target benchmark} & \textbf{Model} & \textbf{Baseline} & \textbf{Direct} & \textbf{Transferred} \\
\midrule
\multirow{3}{*}{OlympiadBench} & \multirow{3}{*}{Omni-MATH} & GPT--5.4      & 56.6 & \na & \dvalp{60.3}{3.7} \\
                                & & GPT--5.4-mini & 34.8 & \na & \dvalp{36.6}{1.8} \\
                                & & GPT--5.4-nano & 38.8 & \na & \dvalp{40.1}{1.3} \\
\bottomrule
\end{tabular}%
}
\caption{Transfer of optimized skills across three axes. \textbf{(a)} \emph{Cross-model}: a skill optimized for the source model is deployed on the target model. \textbf{(b)} \emph{Cross-harness}: a skill trained inside the source harness is evaluated inside the target harness, all on GPT--5.5. \textbf{(c)} \emph{Cross-benchmark}: the source benchmark skill is evaluated on the target benchmark across three target models. \textbf{Baseline} is the target's no-skill score, \textbf{Direct} is the in-domain \ourmethod{} score, and \textbf{Transferred} applies the source skill without further optimization. Subscripts show the change over the target baseline. The GPT--5.4{$\to$}GPT--5.4 transferred cells in (a) are marked \na{} because source and target match (i.e.\ no transfer occurs); we still report the GPT--5.4 baseline and in-domain \ourmethod{} score (taken from Table~\ref{tab:main_results_by_harness}) in those rows for reference. Every remaining row in (a)--(c) is a positive transfer (no row falls below the target's no-skill baseline).}
\label{tab:transfer_all}
\label{tab:transfer_cross_model}
\label{tab:transfer_cross_harness}
\label{tab:transfer_cross_benchmark}
\end{table*}

% --- Legacy: original three separate transfer tables, kept for archival.
% \begin{table*}[t]
% \centering
% \scriptsize
% \setlength{\tabcolsep}{3.5pt}
% \renewcommand{\arraystretch}{0.95}
% \begin{tabular}{l l c c c}
% \toprule
% \textbf{Source skill} & \textbf{Target setting} & \textbf{Baseline} & \textbf{Target \ourmethod{}} & \textbf{Transferred} \\
% \midrule
% GPT--5.4 / SpreadsheetBench         & GPT--5.4 / SpreadsheetBench      & 36.8 & 62.5 & \dvalp{52.1}{15.3} \\
% GPT--5.4 / SpreadsheetBench         & GPT--5.4-mini / SpreadsheetBench & 33.0 & 47.5 & \dvalp{45.5}{12.5} \\
% GPT--5.4 / SpreadsheetBench         & GPT--5.4-nano / SpreadsheetBench & 23.5 & 42.5 & \dvalp{26.5}{3.0}  \\
% GPT--5.4 / LiveMath                 & GPT--5.4 / LiveMath              & 33.6 & 44.0 & \dvalp{47.2}{13.6} \\
% GPT--5.4 / LiveMath                 & GPT--5.4-mini / LiveMath         & 14.7 & 32.8 & \dvalp{19.2}{4.5}  \\
% GPT--5.4 / LiveMath                 & GPT--5.4-nano / LiveMath         & 13.6 & 27.2 & \dvalp{28.8}{15.2} \\
% % GPT--5.4 optimizer / SpreadsheetBench & GPT--5.4-mini / SpreadsheetBench & 36.1 & 47.5 & \dvalp{43.2}{7.1}  \\
% % GPT--5.4 optimizer / SpreadsheetBench & GPT--5.4-nano / SpreadsheetBench & 26.4 & 42.5 & \dvalp{36.4}{10.0} \\
% % GPT--5.4 optimizer / LiveMath       & GPT--5.4-mini / LiveMath         & 24.0 & 32.8 & \dvaln{17.6}{6.4} \\
% % GPT--5.4 optimizer / LiveMath       & GPT--5.4-nano / LiveMath         & 25.6 & 27.2 & \dvalp{26.4}{0.8}  \\
% \bottomrule
% \end{tabular}
% \caption{Cross-model transfer. Each row deploys a source skill onto a different target model (same benchmark). \textbf{Baseline} is the target model's no-skill performance on the target benchmark when available; \textbf{Target \ourmethod{}} is an in-domain reference where a skill is optimized directly for the target; \textbf{Transferred} applies the source skill to the target without further optimization. Gains over the target baseline are shown as small green or red subscripts. For the ALFWorld static rows, we report only the LLM initial skill score and leave the baseline blank.}
% \label{tab:transfer_cross_model_legacy}
% \end{table*}
%
% \begin{table}[t]
% \centering
% \scriptsize
% \setlength{\tabcolsep}{5pt}
% \renewcommand{\arraystretch}{0.95}
% \begin{tabular}{l l l c c c}
% \toprule
% \textbf{Benchmark} & \textbf{Source harness} & \textbf{Target harness} & \textbf{Baseline} & \textbf{Target \ourmethod{}} & \textbf{Transferred} \\
% \midrule
% LiveMath         & Codex       & Claude Code & 40.8 & 49.6 & \dvalp{42.4}{1.6}  \\
% SpreadsheetBench & Codex       & Claude Code & 50.0 & 80.4 & \dvalp{81.8}{31.8} \\
% SpreadsheetBench & Claude Code & Codex       & 41.4 & 85.0 & \dvalp{71.1}{29.7} \\
% \bottomrule
% \end{tabular}
% \caption{Cross-harness transfer (all on GPT--5.5). A skill trained inside one execution harness is evaluated inside another (or the same) harness without re-optimization. Same column semantics as Table~\ref{tab:transfer_cross_model}; gains over the target baseline are shown as a small green subscript.}
% \label{tab:transfer_cross_harness_legacy}
% \end{table}
%
% \begin{table}[t]
% \centering
% \scriptsize
% \setlength{\tabcolsep}{5pt}
% \renewcommand{\arraystretch}{0.95}
% \begin{tabular}{l l l c c}
% \toprule
% \textbf{Source benchmark} & \textbf{Target benchmark} & \textbf{Model} & \textbf{Baseline} & \textbf{Transferred (raw / aligned)} \\
% \midrule
% % LiveMath          & Omni-MATH     & GPT--5.5      & 86.5 & \dvalp{87.5}{1.0} / \dvalp{87.9}{1.4} \\
% % LiveMath          & OlympiadBench & GPT--5.5      & 60.6 & \dvalp{62.5}{1.9} / \dvaln{60.4}{0.2} \\
% LiveMath          & OlympiadBench & GPT--5.4      & 62.7 & \dvaln{61.9}{0.8} / \dvalp{64.8}{2.1} \\
% % OlympiadBench     & Omni-MATH     & GPT--5.4      & 57.6 & \dvaln{56.2}{1.4} / \dvalp{58.7}{1.1} \\
% OlympiadBench     & Omni-MATH     & GPT--5.4      & 56.6 & \dvalp{59.3}{2.7} / \dvalp{60.3}{3.7} \\
% OlympiadBench     & Omni-MATH     & GPT--5.4-mini & 34.8 & \dvalp{36.6}{1.8} / \dvalp{35.6}{0.8} \\
% OlympiadBench     & Omni-MATH     & GPT--5.4-nano & 38.8 & \dvalp{40.1}{1.3} / \dvaln{35.0}{3.8} \\
% \bottomrule
% \end{tabular}
% \caption{Cross-benchmark transfer on math. The source benchmark's optimized skill is evaluated on a related target math benchmark using the same model. \textbf{Baseline} is the target model's no-skill performance on the target benchmark; \textbf{Transferred} reports raw transfer / target-aligned transfer, with gains over baseline shown as small green or red subscripts.}
% \label{tab:transfer_cross_benchmark_legacy}
% \end{table}
